
\section{Related Work}
\label{sec:related-work}

\subsection{Session Types}
\label{sec:session-types}

Context-free session types (CFSTs) have been introduced by
\citet{DBLP:conf/icfp/ThiemannV16} and further developed by
\citet{DBLP:journals/toplas/Padovani19} and
\citet{DBLP:journals/iandc/AlmeidaMTV22}. CFSTs liberate traditional 
session types from their tail-bias by replacing the prefix type
operators for the session operations by symmetric composition. This
move complicates the metatheory as the syntax of types is now
quotiented over the monoid laws and distributivity of composition over
choice. This complication makes type equivalence a decidable, but
difficult problem \cite{DBLP:conf/tacas/AlmeidaMV20} and subtyping an
undecidable problem
\cite{DBLP:journals/toplas/Padovani19,DBLP:journals/tcs/SilvaMV26}.

CFSTs are one starting point of our investigation. Symmetric
composition makes them well-suited to splitting local borrows. Only
remote borrows require extra synchronization machinery. Compared to
the CFST work and the implemented system FreeST based on it
\cite{almeida26:_frees_concur_progr_languag_with}, our calculus
currently does not support polymorphism. We rely on the FreeST
implementation as an external constraint solver for splitting
constraints. 


The most closely related work is ``Borrowing from Session Types'' (BST) \cite{DBLP:journals/pacmpl/SaffrichS0V25}.
They present a similar approach to borrowing, but for regular session types, rather than CFST.
Our type system is heavily inspired from theirs, as it also relies on
BI-contexts that enforce a run-time ordering on operations.
We do rely on their translation from the borrowing notation to
explicit splits, but extend it to cater for the difference between
local and remote splits.
However, our work differs in a number of significant aspects.
We present a direct operational semantics for \thecalculus,
whereas their system's semantics is defined by translation.
We develop a new, efficient representation for borrows, which
distinguishes between local and remote borrows, whereas their
borrows are all remote, i.e., synchronize always. Our translation from \thecalculus to a low
level calculus serves as a blueprint for an implementation and we
present a simulation result.
Metatheoretical results are mechanized using the Agda proof
assistant~\cite{Agda}.

Here are the novel aspects of \thecalculus in direct comparison to
BST.
\begin{itemize}
\item Introduced binder groups to construct the direct semantic model.
\item Introduced the concepts of mobility and boundedness
  (\cref{fig:bounded-session-types}), which are key in enabling
  efficient local borrows.
\item Simplified the typing rules for let expressions and pair
  eliminators: in \thecalculus, they draw their assumption only from a left context,
  which simplifies the treatment of effects and the implementation; we
  observed no impact on example programs.
\item Developed algorithmic typing based on constraints: the BST work
  provided a direct splitting algorithm for borrows based on
  traditional, tail-recursive session types; we leave the development
  of such a dedicated algorithm for future work.
\item Our approach to constraint solving is new.
\item Developed a new, low-level target calculus along with a
  translation into it.
\end{itemize}


%% other notions like Pfenning's et al stuff?
Manifest sharing \cite{DBLP:journals/pacmpl/BalzerP17} and
manifest deadlock freedom \cite{DBLP:conf/esop/BalzerTP19} (an
extension of the former work) is concerned with session types, which
make part of a protocol sharable. 



%% Nested session types \cite{DBLP:journals/toplas/DasDMP22}

\subsection{Ordered Types}
\label{sec:ordered-types}

Law and order \cite{DBLP:journals/pacmpl/Saffrich0024} is a precursor
of BST. It draws inspiration from work on
non-commutative logic and ordered lambda calculus
\cite{DBLP:conf/tlca/PolakowP99,
  DBLP:journals/entcs/PolakowP99}
and puts it to work in the context of typestate. It targets sequential
programming, so it is not concerned with synchronization issues, at all.

Other work on ordered types and lambda calculi is concerned with
memory management and memory layout
\cite{DBLP:conf/popl/PetersenHCP03} or stack allocation
\cite{DBLP:conf/tldi/AhmedW03}. Like Law and order, BST, and the
present work, these earlier works deviate quickly from the pure,
proof-theoretic form of ordered lambda calculus \cite{DBLP:conf/tlca/PolakowP99,
  DBLP:journals/entcs/PolakowP99}. Each develops a calculus that picks
up ideas, but puts significant domain-specific elements on top, to the
extent that even basic elimination principles
change. \citet{DBLP:conf/tldi/AhmedW03} share with our work the use of
BI-contexts to distinguish ordered and unordered binding groups.

\citet{DBLP:journals/pacmpl/CutlerWNHGSP24} present an application of
ordered types to typing a calculus for streaming data. This system is
very close to the original theory of ordered types, but it is
restricted to first-order functions.

\subsection{Borrowing}
\label{sec:borrowing}

Rust

Ahmed and friends


%%% Local Variables:
%%% mode: latex
%%% TeX-master: "../popl2027.tex"
%%% End:
